\documentclass[11pt, a4paper]{article}
\usepackage[utf8]{inputenc}
\usepackage[spanish]{babel}
\usepackage{amsmath, amssymb}
\usepackage{geometry}
\usepackage{tabularx}
\usepackage{booktabs}
\usepackage{enumitem}
\usepackage{xcolor}
\usepackage{fancyhdr}
\usepackage{titlesec}

\geometry{left=2.5cm, right=2.5cm, top=2.5cm, bottom=2.5cm}
\pagestyle{fancy}
\fancyhf{}
\fancyhead[L]{\textbf{Ingeniería de Sistemas Aeronáuticos}}
\fancyhead[R]{\textbf{ATA 21: B737 vs A320}}
\fancyfoot[C]{\thepage}

\titleformat{\section}{\Large\bfseries\color{blue!70!black}}{\thesection}{1em}{}[\titlerule]
\titleformat{\subsection}{\large\bfseries\color{blue!50!black}}{\thesubsection}{1em}{}

\title{\vspace{-2cm}\textbf{GUÍA DE TRABAJO AUTÓNOMO: Análisis Comparativo ATA 21}\\ \large Sistemas de Aire Acondicionado y Presurización (B737 vs A320)}
\author{Unidad 2: Sistemas de Aeronaves \\ \small Tiempo estimado: 60--90 minutos}
\date{}

\begin{document}
\maketitle
\vspace{-1cm}

\noindent\textbf{Nombre:} \rule{6cm}{0.4pt} \hfill \textbf{Matrícula:} \rule{5cm}{0.4pt} \\
\vspace{0.2cm}
\noindent\textbf{Fecha:} \rule{4cm}{0.4pt}

\section*{Objetivo}
Comparar el sistema ATA 21 (aire acondicionado y presurización) entre Boeing 737 y Airbus A320 usando tablas y cálculos básicos.

\section*{Entregables (todo en un solo documento PDF o Word, respuestas numeradas)}
\begin{enumerate}
    \item Tabla comparativa de temperaturas en el ciclo ACM (completar con datos dados).
    \item Cálculo de caída de temperatura en turbina (B737).
    \item Cálculo de altitud de cabina (A320).
    \item Análisis de falla: ¿qué pasa si falla el \textit{Zone Controller} en A320?
    \item Recomendación operacional (máx.\ 3 líneas).
    \item Tabla comparativa rápida B737 vs A320.
\end{enumerate}

\section{Paso 1: Tabla de temperaturas del ACM (B737)}
\textit{Datos proporcionados (no requiere buscar en manual):}

\begin{table}[h]
\centering
\caption{Ciclo ACM del Boeing 737}
\begin{tabular}{|c|l|c|c|}
\hline
\textbf{Etapa} & \textbf{Componente} & \textbf{T entrada (°C)} & \textbf{T salida (°C)} \\ \hline
1 & Primary Heat Exchanger & 215.6 & 76.7 \\ \hline
2 & Compresor ACM & 76.7 & 110 \\ \hline
3 & Secondary Heat Exchanger & 110 & 65.6 \\ \hline
4 & Turbina ACM & 65.6 & 1.6 \\ \hline
\end{tabular}
\end{table}

\textbf{Pregunta:} ¿Cuánto se enfría el aire en la turbina? \\[2pt]
\[
\Delta T = 65.6 - 1.6 = \boxed{64.0 \;^{\circ}\text{C}}
\]
Explica con tus palabras por qué la turbina es la etapa que más enfría.

\vspace{0.3cm}
\textbf{Espacio para respuesta:}
\vspace{1cm}

\section{Paso 2: Cálculo isentrópico de la turbina (B737)}
Fórmula:
\[
T_{out} = T_{in} \times \left( \frac{P_{out}}{P_{in}} \right)^{\frac{\gamma-1}{\gamma}}
\]

\textbf{Datos:}
\begin{itemize}[noitemsep]
    \item \(T_{in} = 65.6 \;^{\circ}\text{C} = 338.75 \text{ K}\)
    \item \(P_{in} = 45 \text{ psi},\quad P_{out} = 14.7 \text{ psi}\)
    \item \(\gamma = 1.4\)
\end{itemize}

\textbf{Realiza los pasos:}
\begin{enumerate}
    \item Relación de presiones: \( \dfrac{14.7}{45} = \boxed{0.3267} \)
    \item Exponente: \( \dfrac{1.4 - 1}{1.4} = \boxed{0.2857} \)
    \item \( T_{out} = 338.75 \times (0.3267)^{0.2857} \approx \boxed{246.5 \text{ K}} \)
    \item En °C: \( 246.5 - 273.15 = \boxed{-26.65 \;^{\circ}\text{C}} \) \\
          En °F: \( (-26.65 \times 9/5) + 32 = \boxed{-15.97 \;^{\circ}\text{F}} \)
\end{enumerate}

\textbf{Compara} con el valor real del manual (1.6\,°C). ¿Por qué existe diferencia?

\vspace{0.3cm}
\textbf{Espacio para respuesta:}
\vspace{1.5cm}

\section{Paso 3: Presurización -- Altitud de cabina (A320)}
\textbf{Datos:}
\begin{itemize}[noitemsep]
    \item Vuelo a FL410 $\rightarrow$ presión exterior $\approx$ 2.72 psi
    \item \(\Delta P_{\max}\) A320 = 8.06 psi
    \item Presión nivel del mar = 14.696 psi
\end{itemize}

\textbf{Fórmula de altitud:}
\[
\text{Altitud (ft)} = 145442 \times \left[ 1 - \left( \frac{P_{cab}}{14.696} \right)^{0.1902} \right]
\]

\textbf{Calcula:}
\begin{enumerate}
    \item \( P_{cab} = 2.72 + 8.06 = \boxed{10.78 \text{ psi}} \)
    \item \( \dfrac{P_{cab}}{14.696} = \dfrac{10.78}{14.696} = \boxed{0.7336} \)
    \item Eleva a 0.1902: \( 0.7336^{0.1902} \approx \boxed{0.9405} \)
    \item \( 1 - 0.9405 = \boxed{0.0595} \)
    \item Multiplica: \( 145442 \times 0.0595 \approx \boxed{8654 \text{ ft}} \)
\end{enumerate}

\textbf{Pregunta:} Según FAR 25, la altitud de cabina no debe exceder 8000 ft en condiciones normales. ¿Este vuelo cumple? ¿Qué implicaciones tiene?

\vspace{0.3cm}
\textbf{Espacio para respuesta:}
\vspace{1.5cm}

\section{Paso 4: Escenario de falla (Airbus A320)}
\textbf{Situación:} En crucero, el \textit{Zone Controller} falla completamente. \\
Según el manual Lufthansa A320 (dato proporcionado): el \textit{Pack Controller} entrega una temperatura fija de \textbf{15\,°C} como fallback.

\textbf{Preguntas:}
\begin{enumerate}
    \item Si la tripulación había seleccionado 24\,°C para la cabina, ¿cuánto se desvía la temperatura real?
    \item ¿Qué medidas puede tomar la tripulación? (descender, aumentar flujo, etc.)
    \item Escribe una \textbf{recomendación operacional} en máximo 3 líneas.
\end{enumerate}

\vspace{0.3cm}
\textbf{Espacio para respuestas:}
\vspace{3cm}

\section{Paso 5: Tabla comparativa rápida}
Completa con \checkmark\ / \textsf{X} o texto breve:

\begin{table}[h]
\centering
\caption{Comparación B737 vs A320}
\begin{tabularx}{\textwidth}{l X X}
\toprule
\textbf{Característica} & \textbf{Boeing 737} & \textbf{Airbus A320} \\ \midrule
Usa ACM de 2 ruedas & & \\ \midrule
Temperatura típica a la salida del pack & $\sim$35\,°F (1.6\,°C) & $\sim$2 a 70\,°C (variable) \\ \midrule
Control de temperatura por zona & Trim Air & Zone Controller + Trim Air \\ \midrule
Altitud máxima de cabina a FL410 (con $\Delta P_{\max}$) & $\sim$8000 ft (NG) & $\sim$8650 ft \\ \bottomrule
\end{tabularx}
\end{table}

\vspace{0.5cm}
\noindent\rule{\textwidth}{0.4pt}

\section*{Rúbrica de Evaluación (Total: 30 puntos)}
\begin{table}[h]
\centering
\begin{tabularx}{\textwidth}{l X}
\toprule
\textbf{Criterio} & \textbf{Puntaje} \\ \midrule
Paso 1: Tabla de temperaturas completa & 5 \\ \midrule
Paso 2: Cálculo de turbina (procedimiento + resultado) & 10 \\ \midrule
Paso 3: Cálculo de altitud de cabina (pasos + interpretación) & 8 \\ \midrule
Pasos 4 y 5: Análisis de falla + recomendación + tabla comparativa & 7 \\ \bottomrule
\textbf{Total} & \textbf{30} \\ \bottomrule
\end{tabularx}
\end{table}

\vspace{0.3cm}
\noindent\textbf{Formato de entrega:} PDF o Word, con nombre, fecha, respuestas numeradas. \\
\noindent\textbf{No se requiere LaTeX.}

\end{document}
